% 中文摘要頁
\begin{ZhAbstract}
    \begin{ZhAbstractItems}
        \noindent \text 關鍵詞：大型語言模型、多代理系統、自動化軟體測試、RESTful API、強化學習、模型上下文協定（MCP）、代理人對代理人協定（A2A）、工作流程決策、軟體品質保證
    \end{ZhAbstractItems}

    \begin{ZhAbstractDescription}
        
        \text 隨著軟體開發邁向敏捷與微服務架構，RESTful API已成為系統整合與服務交換之核心介面。然而，現有自動化測試方法在面對多步驟操作依賴、語意限制條件與真實執行環境變動時，仍常出現語意理解不足、跨操作依賴推理困難，以及測試流程缺乏恢復能力等問題，導致無效請求比例偏高，並限制整體測試可完成度。
        \text 為因應上述問題，本研究提出一套以後執行工作流程恢復層（Post-Execution Workflow Recovery Layer）為核心之多代理RESTful API自動化測試框架。此框架以大型語言模型（Large Language Model, LLM）為語意推理核心，整合AutoGen之代理協作能力、LangGraph之工作流程控制機制，以及模型上下文協定（Model Context Protocol, MCP）與代理人對代理人協定（Agent-to-Agent Protocol, A2A），建立由OpenAPI規格解析、測試案例生成、測試程式產生、API執行、工作流程決策到結果分析之完整測試流程。
        \text 本研究之重點不在於單純產生測試請求，而在於處理API執行後之工作流程決策問題。為此，本文將執行後狀態轉移形式化為馬可夫決策過程（Markov Decision Process, MDP），使系統能根據HTTP狀態碼、錯誤類型、依賴狀態、授權狀態與資源可用性等訊號，選擇後續工作流程行動，例如重試、更新授權、資源解析、重新產生測試資料或進入結果回報，以提升測試流程之持續執行能力與恢復能力。其中，工作流程恢復層負責在404、資源缺失與依賴失敗等情境下提供可持續執行之恢復機制；工作流程決策策略則負責在候選恢復行動中選擇下一步動作。
        \text 在實證評估方面，本研究以TMDB資料集之100筆manual-doc-grounded任務作為主實驗基準，比較v1至v5不同版本之測試表現。結果顯示，v1至v5之通過率分別為0.25、0.72、0.69、0.69與0.79。其中，v5之Strict Pass Rate仍為0.69，與v3及v4相同，但其Final Pass Rate提升至0.79，且新增10筆Workflow Recovered Pass，顯示其主要效益來自執行後工作流程恢復，而非前段測試生成能力之提升。進一步分析顯示，目前恢復增益主要集中於404資源依賴缺失情境。
        \text 在離線強化學習方面，本研究利用工作流程轉移事件建立離線資料集，並以表格式擬合Q值迭代訓練Offline-Q Policy。此部分之主要目的，在於驗證工作流程決策問題是否可被形式化為learning-based workflow decision strategy之可行性探索。結果顯示，學習型策略已能掌握部分工作流程決策模式，並於特定恢復情境中提供額外效益；然而，受限於單一API領域、Live API環境波動與離線資料規模有限，其整體表現尚不足以宣稱全面超越rule-based baseline。整體而言，本研究驗證了大型語言模型、多代理協作與工作流程決策策略整合於RESTful API自動化測試之可行性，並為後續智慧化API測試系統之設計與評估提供具體參考。

    \end{ZhAbstractDescription}
    
\end{ZhAbstract}